\documentclass[10pt,a4paper]{amsart}
\usepackage{amsmath,amssymb,amsthm,graphicx}
\usepackage{verbatim}
\newenvironment{reference}{\comment}{\endcomment}
%\newenvironment{reference}{}{}
\newenvironment{slogan}{\comment}{\endcomment}
\newenvironment{history}{\comment}{\endcomment}

% For commutative diagrams we use Xy-pic
\usepackage[all]{xy}

% We use 2cell for 2-commutative diagrams.
\xyoption{2cell}
% Theorem environments.
%
\theoremstyle{plain}
\newtheorem{theorem}[subsection]{Theorem}
\newtheorem{proposition}[subsection]{Proposition}
\newtheorem{lemma}[subsection]{Lemma}

\theoremstyle{definition}
\newtheorem{definition}[subsection]{Definition}
\newtheorem{example}[subsection]{Example}
\newtheorem{exercise}[subsection]{Exercise}
\newtheorem{situation}[subsection]{Situation}

\theoremstyle{remark}
\newtheorem{remark}[subsection]{Remark}
\newtheorem{remarks}[subsection]{Remarks}

\numberwithin{equation}{subsection}

% Macros
%
\def\lim{\mathop{\mathrm{lim}}\nolimits}
\def\colim{\mathop{\mathrm{colim}}\nolimits}
\def\Spec{\mathop{\mathrm{Spec}}}
\def\Hom{\mathop{\mathrm{Hom}}\nolimits}
\def\Ext{\mathop{\mathrm{Ext}}\nolimits}
\def\SheafHom{\mathop{\mathcal{H}\!\mathit{om}}\nolimits}
\def\SheafExt{\mathop{\mathcal{E}\!\mathit{xt}}\nolimits}
\def\Sch{\mathit{Sch}}
\def\Mor{\mathop{\mathrm{Mor}}\nolimits}
\def\Ob{\mathop{\mathrm{Ob}}\nolimits}
\def\Sh{\mathop{\mathit{Sh}}\nolimits}
\def\NL{\mathop{N\!L}\nolimits}
\def\CH{\mathop{\mathrm{CH}}\nolimits}
\def\proetale{{pro\text{-}\acute{e}tale}}
\def\etale{{\acute{e}tale}}
\def\QCoh{\mathit{QCoh}}
\def\Ker{\mathop{\mathrm{Ker}}}
\def\Im{\mathop{\mathrm{Im}}}
\def\Coker{\mathop{\mathrm{Coker}}}
\def\Coim{\mathop{\mathrm{Coim}}}

% Boxtimes
%
\DeclareMathSymbol{\boxtimes}{\mathbin}{AMSa}{"02}

%
% Macros for moduli stacks/spaces
%
\def\QCohstack{\mathcal{QC}\!\mathit{oh}}
\def\Cohstack{\mathcal{C}\!\mathit{oh}}
\def\Spacesstack{\mathcal{S}\!\mathit{paces}}
\def\Quotfunctor{\mathrm{Quot}}
\def\Hilbfunctor{\mathrm{Hilb}}
\def\Curvesstack{\mathcal{C}\!\mathit{urves}}
\def\Polarizedstack{\mathcal{P}\!\mathit{olarized}}
\def\Complexesstack{\mathcal{C}\!\mathit{omplexes}}
% \Pic is the operator that assigns to X its picard group, usage \Pic(X)
% \Picardstack_{X/B} denotes the Picard stack of X over B
% \Picardfunctor_{X/B} denotes the Picard functor of X over B
\def\Pic{\mathop{\mathrm{Pic}}\nolimits}
\def\Picardstack{\mathcal{P}\!\mathit{ic}}
\def\Picardfunctor{\mathrm{Pic}}
\def\Deformationcategory{\mathcal{D}\!\mathit{ef}}
\usepackage{cleveref}
\crefname{originalreference}{original reference}{original references}
\Crefname{originalreference}{Original reference}{Original references}
\makeatletter
\newlabel{lemma-H-0-separable}{{0FI0}{source}{}{sourcepointer-1}{}}
\newlabel{lemma-H-0-separable@cref}{{[originalreference][0][]0FI0}{[1][0][]source}{}{}{}}
\makeatother
\crefname{originalreference}{original reference}{original references}
\Crefname{originalreference}{Original reference}{Original references}
\makeatletter
\newlabel{lemma-weil-additive}{{0FHJ}{source}{}{sourcepointer-2}{}}
\newlabel{lemma-weil-additive@cref}{{[originalreference][0][]0FHJ}{[1][0][]source}{}{}{}}
\makeatother
\title{Negative cohomology produces a failure of the separability conditions}
\author{The Stacks Project Authors}
\date{stacks-project:weil.tex:local-snapshot-20261012; GNU FDL 1.2 or later}
\begin{document}
\maketitle
\noindent\textit{Editorial note.} All delivered mathematical text is the author's own, copied from the pinned Stacks Project chapter file weil.tex; no statement or proof step was rewritten, reordered or removed.\par\medskip
\noindent\textit{Editorial note.} The counted proof is the author's complete argument: the reduction to an even negative degree, the multiplication by a power of the projective line, and the direct summand argument showing the failure of the characterizations.\par\medskip
\noindent\textit{Editorial note.} The label lemma-negative-cohomology names the counted statement, and the proof environment immediately following it in the source is delivered verbatim as its proof.\par\medskip
\noindent\textit{Editorial note.} The counted statement is carried by the author's own lemma environment in the chapter source.\par\medskip
\noindent\textit{Editorial note.} Pre-open gate: python operations/agent-authoring/tools/precard\_gate.py --source weil.tex --label lemma-negative-cohomology --start 3928 --end 3951 --local-id 28159 returned RESULT=OK (L28159 free, occupancy FREE, public interval clear). Pre-assembly tag self-check: canonical tag resolved and not present in reference/public-items.jsonl.\par\medskip
\noindent\textit{Editorial note.} The statement and proof are copied from the pinned chapter file; the cited inputs are kept as external pointers with their canonical Stacks tags.\par\medskip
\noindent\textit{Allowed standard prerequisites (editorial).}\begin{itemize}
\item Additivity of the theory under disjoint unions. Same chapter, the lemma computing the cohomology of a disjoint union. Reduces to an irreducible, hence equidimensional, example.
\item The cohomology of the projective line. Same chapter, the computation of the cohomology of the projective line, one dimensional in degrees zero and two. Provides a nonzero class in degree two used to shift a negative degree into an even negative degree.
\item Characterizations of the degree zero cohomology. Same chapter, the characterizations of the degree zero cohomology through injectivity of the pullbacks to points and through generation of the top cohomology by classes of points. Identifies what it means for the conditions to fail and shows that the constructed direct summand prevents the comparison map from being an isomorphism.
\end{itemize}\medskip
\ifdefined\hypertarget\hypertarget{sourcepointer-1}{}\fi
\noindent\textit{Source reference.} Original reference 0FI0: sources/stacks-weil/pinned-original-source.tex, source lines 3860--3860. Prerequisite pointer invoked at proof line(s) 3934; it names a standard fact of the same chapter that the counted argument uses, and every step of the counted proof built on it is retained verbatim in the delivered body.\par\medskip
\ifdefined\hypertarget\hypertarget{sourcepointer-2}{}\fi
\noindent\textit{Source reference.} Original reference 0FHJ: sources/stacks-weil/pinned-original-source.tex, source lines 2932--2932. Prerequisite pointer invoked at proof line(s) 3938; it names a standard fact of the same chapter that the counted argument uses, and every step of the counted proof built on it is retained verbatim in the delivered body.\par\medskip
\begin{lemma}
\label{lemma-negative-cohomology}
Assume given (D0), (D1), (D2), and (D3) satisfying (A), (B), and (C).
If there exists a smooth projective scheme $Y$ over $k$ such that
$H^i(Y)$ is nonzero for some $i < 0$, then there exists an equidimensional
smooth projective scheme $X$ over $k$ such that the equivalent conditions
of Lemma \ref{lemma-H-0-separable} fail for $X$.
\end{lemma}

\begin{proof}
By Lemma \ref{lemma-weil-additive} we may assume $Y$ is irreducible
and a fortiori equidimensional. If $i$ is odd, then after replacing
$Y$ by $Y \times Y$ we find an example where $Y$ is equidimensional
and $i = -2l$ for some $l > 0$. Set $X = Y \times (\mathbf{P}^1_k)^l$.
Using axiom (B)(a) we obtain
$$
H^0(X) \supset H^0(Y) \oplus
H^i(Y) \otimes_F H^2(\mathbf{P}^1_k)^{\otimes_F l}
$$
with both summands nonzero. Thus it is clear that $H^0(X)$ cannot be
isomorphic to $H^0$ of the spectrum of
$\Gamma(X, \mathcal{O}_X) = \Gamma(Y, \mathcal{O}_Y)$
as this falls into the first summand.
\end{proof}

\end{document}
